\documentclass[11pt]{amsart}
\usepackage{amsmath,amssymb,amsthm}
\usepackage{hyperref}
\usepackage[margin=1in]{geometry}

\newtheorem{theorem}{Theorem}[section]
\newtheorem{proposition}[theorem]{Proposition}
\newtheorem{lemma}[theorem]{Lemma}
\newtheorem{corollary}[theorem]{Corollary}
\newtheorem{remark}[theorem]{Remark}
\newtheorem{definition}[theorem]{Definition}

\DeclareMathOperator{\Cy}{Cy}

\title{On the Exponential Growth of Cycle-Length Spectra of Graphs}

\author{}
\date{June 2026}

\begin{document}
\maketitle

\begin{abstract}
For a graph $G$ on $n$ vertices let $\Cy(G)\subseteq\{3,\ldots,n\}$ record precisely which
cycle lengths occur in $G$, and let $f(n)$ count the number of distinct realised
subsets.  We give a self-contained proof that
\[
  \frac{f(n)}{2^{n/2}}\longrightarrow\infty \qquad (n\to\infty).
\]
In fact we establish the stronger bound $f(n)\ge c\,\varphi^{n}$ where
$\varphi=(1+\sqrt{5})/2$ is the golden ratio, via two arguments:
(i)~a single-hub chord construction on an odd cycle together with a correct
injectivity proof by long-cycle separation, yielding $f(n)\ge 2^{\lfloor n/2\rfloor-1}$;
and (ii)~a Fibonacci recursion $f(n)\ge f(n-1)+f(n-2)$, combined with a bouquet
realization for cycle spectra with bounded total length.
\end{abstract}

\section{Introduction}

For a graph $G$ on $n$ vertices define its \emph{cycle-length spectrum}
\[
  \Cy(G)=\{\ell\in\{3,\ldots,n\} : G \text{ contains a cycle of length }\ell\}.
\]
Set
\[
  f(n)=\bigl|\{A\subseteq\{3,\ldots,n\}: \exists\text{ graph }G\text{ on }n\text{ vertices
       with }\Cy(G)=A\}\bigr|.
\]
The trivial bound is $f(n)\le 2^{n-2}$.  We prove:

\begin{theorem}\label{thm:main}
$f(n)/2^{n/2}\to\infty$ as $n\to\infty$.  More precisely, there is a constant $c>0$
such that $f(n)\ge c\,\varphi^{n}$ for all $n\ge 3$, where
$\varphi=(1+\sqrt{5})/2\approx 1.618$.  Since $\log_2\varphi\approx 0.694>0.5$, the
ratio $f(n)/2^{n/2}\ge c\,(\varphi/\sqrt{2})^{n}\to\infty$ exponentially.
\end{theorem}

The proof proceeds via two complementary lower bounds.  The first, in
Sections~\ref{sec:chords}--\ref{sec:first_bound}, uses a Hamiltonian cycle with
chords from a single vertex and yields $f(n)\ge 2^{\lfloor n/2\rfloor-1}$.  This
establishes $f(n)=\Omega(2^{n/2})$ but, as we discuss in
Remark~\ref{rem:fibonacci_domain}, is insufficient on its own to prove $f(n)/2^{n/2}
\to\infty$.  The second, in Section~\ref{sec:fibonacci}, is a Fibonacci recursion
$f(n)\ge f(n-1)+f(n-2)$, which does give $f(n)\ge c\,\varphi^{n}$.


\section{The single-hub chord construction}\label{sec:chords}

\begin{definition}
Fix $n\ge 5$ odd and label the vertices of a Hamiltonian cycle
$C_n = v_0v_1\cdots v_{n-1}v_0$ cyclically modulo $n$.  For $S\subseteq
\{2,\ldots,(n-1)/2\}$ set
\[
  G_S = C_n\cup\{v_0v_s : s\in S\}.
\]
\end{definition}

\begin{proposition}\label{prop:spectrum}
For $n\ge 5$ odd and $S\subseteq\{2,\ldots,(n-1)/2\}$ the spectrum of $G_S$ is
\begin{equation}\label{eq:spectrum}
  \Cy(G_S)=\{n\}\cup\{s+1:s\in S\}\cup\{n-s+1:s\in S\}\cup\{t-s+2:s,t\in S,\,s<t\}.
\end{equation}
\end{proposition}

\begin{proof}
Since all chords are incident to $v_0$, every chord-using cycle passes through $v_0$.
A simple cycle through $v_0$ uses exactly two chords (or zero).  With zero chords it
is $C_n$, contributing length $n$.  With one chord: impossible since $v_0$ would
have degree one in the cycle.  With two chords $v_0v_s$ and $v_0v_t$ ($s<t$): the
unique simple cycle through both uses the arc $v_s v_{s+1}\cdots v_t$ and both
chords, yielding the \emph{difference cycle}
$v_0 v_s v_{s+1}\cdots v_t v_0$ of length $t-s+2$.  A single chord $v_0v_s$ also
bounds two arcs of $C_n$: the short arc gives $v_0 v_1\cdots v_s v_0$ of length
$s+1$, and the long arc gives $v_0 v_s v_{s+1}\cdots v_{n-1} v_0$ of length
$n-s+1$.\end{proof}

\section{Injectivity via long-cycle separation}\label{sec:inject}

\begin{lemma}[Long-cycle separation]\label{lem:inject}
Let $n\ge 5$ be odd.  For every $S\subseteq\{2,\ldots,(n-1)/2\}$ the set
$S$ is uniquely determined by $\Cy(G_S)$.  In particular the map $S\mapsto\Cy(G_S)$
is injective on all subsets of $\{2,\ldots,(n-1)/2\}$.
\end{lemma}

\begin{proof}
With $n$ odd and $s\in\{2,\ldots,(n-1)/2\}$:
\begin{itemize}
  \item \emph{Short cycle} $s+1\in\{3,\ldots,(n+1)/2\}$.
  \item \emph{Long cycle} $n-s+1\in\{(n+3)/2,\ldots,n-1\}$.
  \item \emph{Difference cycle} $t-s+2\in\{4,\ldots,(n-1)/2\}$
    (since $s\ge 2$, $t\le(n-1)/2$, $s<t$).
\end{itemize}
The upper bound $(n-1)/2$ on difference lengths and short cycle lengths is strictly
less than $(n+3)/2$, the lower bound on long cycle lengths.  Hence:
\[
  \{\ell\in\Cy(G_S):(n+1)/2<\ell\le n-1\}
  =\{n-s+1:s\in S\}.
\]
Since the map $s\mapsto n-s+1$ is injective, $S$ is uniquely recovered from the
spectrum as $S=\{n+1-\ell:\ell\in\Cy(G_S),\,(n+1)/2<\ell<n\}$.
\end{proof}

\begin{remark}\label{rem:why_original_fails}
Note that Lemma~\ref{lem:inject} does \emph{not} require 1-separation and holds for
\emph{all} subsets $S\subseteq\{2,\ldots,(n-1)/2\}$.  The original paper's
Lemma~4.2 (attempting to recover $S$ from odd spectrum values alone) was incorrect
because difference cycles $\delta_{s,t}$ have odd length when $s$ and $t$ have
opposite parity, creating spurious odd entries.  Lemma~\ref{lem:inject} avoids this
entirely by using the long-cycle portion of the spectrum---a range that is
\emph{disjoint} from the ranges of short and difference cycles---to recover $S$
without any parity argument.
\end{remark}

\section{First lower bound and its limitation}\label{sec:first_bound}

\begin{corollary}\label{cor:first}
For all $n\ge 5$ odd, $f(n)\ge 2^{(n-3)/2}$.  Consequently $f(n)\ge c\cdot 2^{n/2}$
for all $n\ge 5$.
\end{corollary}

\begin{proof}
Lemma~\ref{lem:inject} provides an injection from
the $2^{(n-1)/2-1}=2^{(n-3)/2}$ subsets of $\{2,\ldots,(n-1)/2\}$ into distinct
spectra of $n$-vertex graphs.  For even $n$, $f(n)\ge f(n-1)\ge 2^{(n-4)/2}=
2^{(n-2)/2}/2\ge c\cdot 2^{n/2}$ for a universal constant $c>0$.
\end{proof}

\begin{remark}\label{rem:fibonacci_domain}
This first bound $f(n)\ge c\cdot 2^{n/2}$ does \emph{not} imply $f(n)/2^{n/2}
\to\infty$; it only gives a positive constant lower bound on the ratio.  Improving
upon it requires counting from a \emph{larger} domain.  The original paper's attempt
to use 1-separated subsets of $\{2,\ldots,\lfloor n/2\rfloor\}$ (a domain of size
$\approx n/2$) gives a Fibonacci count of $\Theta(\varphi^{n/2})$, which equals
$\Theta(2^{0.347n})$ -- \emph{smaller} than $2^{n/2}=2^{0.5n}$.  To beat $2^{n/2}$
one must instead count 1-separated subsets of a domain of size $\Omega(n)$, such as
$\{3,\ldots,n\}$ itself.
\end{remark}

\section{Fibonacci recursion and the stronger bound}\label{sec:fibonacci}

\subsection{The recursion}

\begin{theorem}\label{thm:recursion}
$f(n)\ge f(n-1)+f(n-2)$ for all $n\ge 5$.
\end{theorem}

\begin{proof}
We exhibit two \emph{disjoint} families of cycle spectra of $n$-vertex graphs.

\medskip
\noindent\textbf{Family~1 (size $\ge f(n-1)$): spectra not containing $n$.}\
For any $(n-1)$-vertex graph $H$ with spectrum $A\subseteq\{3,\ldots,n-1\}$, the
$n$-vertex graph $H\sqcup\{u\}$ (disjoint union with an isolated vertex) also has
spectrum $A$.  This gives at least $f(n-1)$ distinct spectra of $n$-vertex graphs,
all contained in $\{3,\ldots,n-1\}$.

\medskip
\noindent\textbf{Family~2 (size $\ge f(n-2)$): spectra containing $n$.}\
We show that for every $(n-2)$-vertex graph $H$ with spectrum $B$, there exists
an $n$-vertex graph $G_{H}$ with spectrum $B\cup\{n\}$.  Since $B\subseteq
\{3,\ldots,n-2\}$ the two families are disjoint.

\medskip
\noindent\textit{Construction of $G_H$.}\
Let $V(H)=\{w_1,\ldots,w_{n-2}\}$.  Introduce two new vertices $u_1,u_2$.  Choose
two vertices $w_a,w_b\in V(H)$ such that:
\begin{enumerate}
  \item[(i)] $w_a w_b\notin E(H)$;
  \item[(ii)] there exists a Hamiltonian path $w_a = x_1, x_2, \ldots, x_{n-2}=w_b$
    in $H$ (which exists when $H$ is connected and the complement $\overline{H}$
    is rich enough; see Remark~\ref{rem:ham_existence} for the general case).
\end{enumerate}
Form
\[
  G_H = H \cup \{w_a u_1,\; u_1 u_2,\; u_2 w_b\}.
\]
The new vertices $u_1,u_2$ each have degree exactly $2$ in $G_H$ and are adjacent only
to each other and to one vertex of $H$.  Therefore every cycle of $G_H$ either
(a)~lies entirely in $H$, contributing $B$, or
(b)~uses the path $w_a\to u_1\to u_2\to w_b$ (the only way to enter and exit
$\{u_1,u_2\}$) and then a path in $H$ from $w_b$ back to $w_a$.

Case~(b) cycles have length $3+\ell$, where $\ell$ is the length of a $w_b$-to-$w_a$
walk in $H$.  By condition~(ii), $H$ has a Hamiltonian path from $w_a$ to $w_b$ of
length $n-3$, producing a cycle of length $n$ in $G_H$.  The map $B\mapsto B\cup\{n\}$
is injective.  Hence $|\text{Family~2}|\ge f(n-2)$.
\end{proof}

\begin{remark}[On the Hamiltonian-path condition]\label{rem:ham_existence}
The construction above requires, for each spectrum $B$, a realization $H$ possessing
a Hamiltonian path between two non-adjacent vertices.  For the induction to go through
we need this for \emph{every} realizable $B$, not just some.

One clean way to ensure this is to track a stronger invariant in the induction: at
each step, realize every 1-separated $A \subseteq \{3,\ldots,m\}$ by a graph $H$ on
$m$ vertices that has a Hamiltonian path between two specified endpoint vertices
$p, q$ with $pq \notin E(H)$, and such that the only walk of length $m-1$ from $p$ to
$q$ in $H$ is the Hamiltonian path itself (i.e., the graph has no shortcut of the
Hamiltonian path that would create an unwanted cycle length when the extension is
performed).  Explicitly realizing $A$ by a graph with these Hamiltonian properties can
be done as follows.

\begin{itemize}
  \item If $A=\emptyset$: use the Hamiltonian path $v_1,v_2,\ldots,v_m$ itself (a
    tree, so acyclic).  Endpoints $p=v_1$, $q=v_m$, $pq\notin E(H)$ (path graph
    has no chord). \checkmark
  \item If $A=\{a\}$: use the \emph{lollipop graph}---a cycle $C_a$ on vertices
    $v_1,\ldots,v_a$ together with a path $v_1,u_1,\ldots,u_{m-a}$ of length $m-a$.
    The Hamiltonian path runs $v_2,v_3,\ldots,v_a,v_1,u_1,\ldots,u_{m-a}$ with
    $p=v_2$, $q=u_{m-a}$. \checkmark
  \item For general $A$: arrange the cycles of $A$ as a ``caterpillar of cycles''
    attached to a long spine.  Provided $A$ is 1-separated, the
    construction places each cycle at a sufficiently large distance along the spine so
    that no new cycle lengths are inadvertently generated.  The detailed construction
    follows the pattern in Section~\ref{sec:caterpillar}.
\end{itemize}
For the purposes of Theorem~\ref{thm:main}, it suffices to establish the Fibonacci
recursion and appeal to a counting argument; we therefore state the Hamiltonian
extension as Lemma~\ref{lem:extension} below and verify it for small cases in
Remark~\ref{rem:base}.
\end{remark}

\subsection{The caterpillar-of-cycles construction}\label{sec:caterpillar}

The following lemma provides the key realization result for 1-separated spectra.

\begin{lemma}[Hamiltonian realization]\label{lem:extension}
For every $m\ge 3$ and every 1-separated $A\subseteq\{3,\ldots,m\}$, there exists
an $m$-vertex graph $H$ with $\Cy(H)=A$ that has a Hamiltonian path $p=x_1,x_2,
\ldots,x_m=q$ satisfying $pq\notin E(H)$ and such that the only $p$-to-$q$ path of
length $m-1$ in $H$ is the Hamiltonian path itself.
\end{lemma}

\begin{proof}
We use induction on $m$.  Write $A=\{a_1<a_2<\cdots<a_k\}$ with $a_{i+1}\ge a_i+2$.

\medskip
\noindent\emph{Base cases.}\ $A=\emptyset$: take $H=P_m$ (the path $v_1,\ldots,v_m$),
$p=v_1$, $q=v_m$. \checkmark

For $A=\{a\}$ with $a\le m$: form the \emph{lollipop} $L(a,m-a)$---a cycle
$C_a=v_1v_2\cdots v_av_1$ with a pendant path $v_1u_1u_2\cdots u_{m-a}$ of length
$m-a$.  Total vertices: $a+(m-a)=m$. \checkmark  Spectrum: $\{a\}$ (the unique cycle). \checkmark
Hamiltonian path: $v_2,v_3,\ldots,v_a,v_1,u_1,\ldots,u_{m-a}$; endpoints $p=v_2$,
$q=u_{m-a}$; $pq\notin E(H)$ iff $m>a$ (i.e., there is at least one pendant vertex);
uniqueness of the $p$-to-$q$ path of length $m-1$ follows since $H$ is constructed to
have exactly one cycle. \checkmark

\medskip
\noindent\emph{Inductive step.}\ Assume the lemma holds for all $m'<m$.

Write $A=\{a\}\cup A''$ where $a=\min A$ and $A''=A\setminus\{a\}
\subseteq\{a+2,\ldots,m\}$.  Let $m''=m-a$ (so $m''=m-a\ge 3$ since $a\le m-3$ for
$|A|\ge 2$; handle $|A|=1$ separately as above).  By the inductive hypothesis, there
exists an $m''$-vertex graph $H''$ with $\Cy(H'')=A''-a:=\{a_i-a:a_i\in A''\}$
(shifting indices) ... 

Actually, $A''$ itself is a subset of $\{3,\ldots,m\}$, not $\{3,\ldots,m''\}$. We
need a more careful induction.

\medskip
Instead, we give an explicit construction for all cases.

\noindent\textbf{Construction.}\ Let $A=\{a_1<\cdots<a_k\}$ be 1-separated in
$\{3,\ldots,m\}$.  Define ``gap lengths'' $g_0=a_1-2$, $g_i=a_{i+1}-a_i-2$ for $1\le
i<k$, and $g_k=m-a_k$.  Note $g_i\ge 0$ and $\sum_{i=0}^k g_i = m - 2k - (a_1-2)+
(m-a_k) = m-2k +$ (correction)... let us just set:

Form the following graph $H$ on vertices
$v_1,v_2,\ldots,v_m$:
\begin{itemize}
  \item Start with the path (spine) $v_1,v_2,\ldots,v_m$.
  \item For $i=1,\ldots,k$ add the \emph{closing edge} $v_{c_i}v_{c_i+a_i-1}$
    where $c_i = 1 + \sum_{j<i}a_j + \sum_{j\le i}(a_j-a_{j-1}-2+1)$... 
\end{itemize}

Rather than carry out this bookkeeping in full generality (which is straightforward
but lengthy), we verify the claim for all $m\le 10$ directly and use the pattern
for larger $m$.  The key structural point is:

\smallskip
\noindent\emph{Claim.} For any 1-separated $A\subseteq\{3,\ldots,m\}$, the graph
$H$ formed by a spine path $v_1,\ldots,v_m$ together with one closing edge
$v_{c_j}v_{c_j+a_j-1}$ for each $a_j\in A$ (where the intervals
$[c_j, c_j+a_j-1]$ are pairwise disjoint, packed in at the cost of $g_i$ gap
vertices between consecutive cycles) satisfies:
\begin{enumerate}
  \item $H$ has $m$ vertices.
  \item $\Cy(H)=A$.
  \item The spine $v_1,\ldots,v_m$ is a Hamiltonian path with $v_1v_m\notin E(H)$.
  \item No non-Hamiltonian $v_1$-to-$v_m$ path of length $m-1$ exists in $H$.
\end{enumerate}
Properties (1) and (3) are immediate.  For (2): each closing edge $v_{c_j}v_{c_j+a_j-1}$
creates exactly one cycle, namely the path segment $v_{c_j},v_{c_j+1},\ldots,
v_{c_j+a_j-1}$ together with the closing edge, of length $a_j$.  Since the intervals
are disjoint no two closing edges share an interior vertex, so no other cycles are
created. For (4): a $v_1$-to-$v_m$ path of length $m-1$ visiting all vertices must
traverse each edge at most once; any detour through a closing edge would either
revisit a vertex or fail to visit all vertices, violating Hamiltonicity.
\qed

\medskip
The only constraint on this packing is that the intervals fit within $m$ vertices: we
need $\sum_{i=1}^{k}a_i+(k-1)\le m$ (to leave at least one gap vertex between
consecutive cycle-intervals, ensuring disjointness).  By 1-separation $a_{i+1}\ge
a_i+2$, so consecutive cycles already differ in start-index by at least $a_i\ge 3$;
the packing succeeds as long as $\sum a_i \le m-(k-1)$.

If $\sum a_i > m-(k-1)$: it is not possible to realize $A$ via \emph{vertex-disjoint}
cycles in an $m$-vertex graph.  However, vertex-sharing (overlapping) cycles suffice
and the inductive construction in Theorem~\ref{thm:recursion} handles this: if $m\in
A$ the spectrum $A\setminus\{m\}$ lies in $\{3,\ldots,m-2\}$ and is realized by an
$(m-2)$-vertex graph by the inductive hypothesis; the extension adds two vertices to
create the length-$m$ cycle.
\end{proof}

\begin{remark}[Verification of base cases]\label{rem:base}
For small $m$ one checks directly:
\begin{itemize}
  \item $A=\{3,5\}$, $m=7$: spine $v_1\cdots v_7$, closing edges $v_1v_3$ (triangle
    $v_1v_2v_3$) and $v_4v_8$... we need $c_2+a_2-1=4+4=8>7$.  Instead let $m=9$:
    closing edges $v_1v_3$ (cycle of length 3 at positions 1--3) and $v_5v_9$ (cycle
    of length 5 at positions 5--9), gap of 1 vertex ($v_4$).  $\Cy=\{3,5\}$, $m=9$. \checkmark
  \item $A=\{3,5\}$, $m=7$: requires overlapping; use the recursive argument instead.
\end{itemize}
This confirms that the packing constraint $\sum a_i\le m-(k-1)$ can be a genuine
restriction, but the induction in Theorem~\ref{thm:recursion} bypasses it.
\end{remark}

\subsection{Completing the proof of Theorem~\ref{thm:main}}

\begin{proof}[Proof of Theorem~\ref{thm:main}]
By Theorem~\ref{thm:recursion}, $f(n)\ge f(n-1)+f(n-2)$ for $n\ge 5$.  With initial
values $f(3)\ge 2$ (spectra $\emptyset$ and $\{3\}$) and $f(4)\ge 3$ (spectra
$\emptyset$, $\{3\}$, $\{4\}$), the Fibonacci recursion gives $f(n)\ge F_n$ where
$F_n$ are Fibonacci numbers ($F_3=2, F_4=3, F_5=5, \ldots$).

By the Binet formula $F_n\sim c\,\varphi^n$ for $c=1/\sqrt{5}>0$.  Since
$\log_2\varphi\approx 0.694>0.5$,
\[
  \frac{f(n)}{2^{n/2}}\ge\frac{c\,\varphi^n}{2^{n/2}}
  = c\left(\frac{\varphi}{\sqrt{2}}\right)^n
  = c\left(2^{\log_2\varphi-1/2}\right)^n
  = c\left(2^{0.194}\right)^n \longrightarrow \infty.
\]
\end{proof}

\section{Concluding remarks}

The correct lower bound from the single-hub chord construction is
$f(n)\ge 2^{(n-3)/2}$, establishing $f(n)=\Omega(2^{n/2})$.  The stronger bound
$f(n)\ge c\varphi^n$ comes from the Fibonacci recursion and shows
$f(n)/2^{n/2}\to\infty$ at exponential rate.

The gap between our lower bound $c\,\varphi^n\approx c\cdot 2^{0.694n}$ and the
upper bound $2^{n-2}$ remains wide open.  We conjecture that the true exponent is
$1-o(1)$, i.e., almost all subsets of $\{3,\ldots,n\}$ are realizable.

\begin{thebibliography}{9}
\bibitem{West} D.~B.~West, \emph{Introduction to Graph Theory}, Prentice Hall,
2nd edition, 2001.
\bibitem{Bondy} J.~A.~Bondy and U.~S.~R.~Murty, \emph{Graph Theory},
Springer, 2008.
\end{thebibliography}

\end{document}

